% Work/MAT — 1. rocnik — hodina 49; navazuje na 48 (deleni mnohoclenu).
% Podklady: ucebnice/01_Zakladni_poznatky/Algebraicke_vyrazy/Algebraické výrazy.tex,
% oddil Rozklad mnohoclenu; Realisticky.cz 010806 Rozklad mnohoclenu na soucin I (vytykani).
% Metodicka inspirace; konkretni priklady a reseni jsou vlastni.
\documentclass[11pt,a4paper]{article}
% Spolecny styl pro ucitelske pripravy MAT (Work/MAT)
% Verze 1.1, 2026-10-23
\NeedsTeXFormat{LaTeX2e}
\ProvidesPackage{priprava}[2026/10/23 v1.1 Sablona ucitelskych priprav MAT]

% Zakladni typografie a cestina
\RequirePackage[T1]{fontenc}
\RequirePackage[utf8]{inputenc}
\RequirePackage[czech]{babel}
\RequirePackage{lmodern}
\RequirePackage{microtype}
\RequirePackage{amsmath,amssymb}
\RequirePackage{xcolor}
\RequirePackage{enumitem}
\RequirePackage[most]{tcolorbox}
\RequirePackage{tabularx,booktabs,array}
\RequirePackage{fancyhdr}
\RequirePackage{needspace}
\RequirePackage{geometry}
\RequirePackage[hidelinks]{hyperref}

\geometry{a4paper,top=17mm,bottom=18mm,left=23mm,right=23mm}
\definecolor{notegray}{gray}{0.39}
\setlength{\parindent}{0pt}
\setlength{\parskip}{2.5pt plus 1pt minus 0.5pt}
\setlist[itemize]{topsep=2pt,itemsep=1pt,parsep=0pt,leftmargin=1.45em}
\setlist[enumerate]{topsep=2pt,itemsep=2pt,parsep=0pt,leftmargin=1.65em}
\setcounter{secnumdepth}{0}
\renewcommand{\arraystretch}{1.12}

% Zahlavi: vlevo rocnik | globalni tema, vpravo cislo: nazev hodiny
% Zapati zustava zcela prazdne.
\pagestyle{fancy}
\fancyhf{}
\renewcommand{\headrulewidth}{0.35pt}
\renewcommand{\footrulewidth}{0pt}
\setlength{\headheight}{14pt}
\newcommand{\nastavPripravu}[4]{%
  \fancyhead[L]{\small #1\enspace\textbar\enspace #2}%
  \fancyhead[R]{\small #3: #4}%
}

% Minimalisticke cernobile boxy. Kazda definice zvlast.
\tcbset{commonlesson/.style={enhanced,breakable=false,colback=white,colframe=black,
  colbacktitle=white,coltitle=black,boxrule=0.45pt,
  arc=0.8mm,outer arc=0.8mm,left=3mm,right=3mm,top=1.6mm,bottom=1.6mm,
  toptitle=1mm,bottomtitle=0.7mm,fonttitle=\bfseries,
  attach boxed title to top left={xshift=2.5mm,yshift=-1.15mm},
  boxed title style={colback=white,colframe=white,boxrule=0pt,left=1mm,right=1mm}}}
\newtcolorbox{zadaniBox}[1][]{commonlesson,title={Zadání},#1}
\newtcolorbox{prikladBox}[1][]{commonlesson,title={Příklad},#1}
\newtcolorbox{definiceBox}[1][]{commonlesson,title={Definice},#1}

% Spolecny box pro zadani a jeho reseni, bez mezery mezi castmi.
% Pouziti:
%   \begin{prikladReseniBox}
%     Zadani ulohy.
%     \tcblower
%     \begin{align*} ... \end{align*}
%   \end{prikladReseniBox}
% Spodni cast ma jemny sedy podklad a delici linku; lze v ni sazet
% vice radku matematiky. Obe casti maji jen jeden spolecny ramecek.
\tcbset{solutionlesson/.style={
  commonlesson,
  enhanced jigsaw,
  breakable,
  colbacklower=black!3,
  segmentation style={solid,draw=black!25,line width=0.35pt},
  fontlower=\small\color{notegray},
  before lower={\noindent\textbf{Řešení:}\enspace},
  leftlower=3mm,rightlower=3mm
}}
\newtcolorbox{prikladReseniBox}[1][]{solutionlesson,title={Příklad},#1}
\newtcolorbox{zadaniReseniBox}[1][]{solutionlesson,title={Zadání},#1}

% Pomucky pro autora pripravy: sedive poznamky; \reseni pro starsi pripravy
\newcommand{\poznamka}[1]{\par\begingroup\small\color{notegray}\itshape
  \textnormal{\textbf{Poznámka pro učitele.}}\enspace #1\par\endgroup}
\newcommand{\reseni}[1]{\par\begingroup\small\color{notegray}
  \textbf{Očekávané řešení:}\enspace #1\par\endgroup}
\newcommand{\blok}[1]{\needspace{5\baselineskip}\subsection*{#1}}
\newcommand{\vyrok}[1]{\uv{#1}}

\nastavPripravu{1. ročník}{Algebraické výrazy}{49}{Vytýkání a seskupování}
\begin{document}
\poznamka{Jádro přibližně 45 minut: 5 min motivace a obrácení distributivity, 9 min vytýkání v P1--P2, 10 min S1 a kontrola, 8 min seskupování v P3, 10 min S2 a kontrola, 3 min shrnutí. Náročnější seskupování a kombinované rozklady jsou rezerva; domácí úlohy 15--20 minut.}
\blok{Od roznásobování k rozkladu}
Dosud jsme uměli násobit a odstraňovat závorky:
\[3x(x+4)=3x^2+12x.\]
Někdy potřebujeme přesný opak: \emph{vytvořit součin}, například kvůli pozdějšímu krácení výrazů nebo řešení rovnic:
\[3x^2+12x=3x(x+4).\]
Obě podoby mají pro všechna $x\in\mathbb R$ stejnou hodnotu.
\begin{definiceBox}
\textbf{Rozklad mnohočlenu na součin} je úprava, při které mnohočlen vyjádříme jako součin jednodušších mnohočlenů se stejnou hodnotou pro všechna dosazení.
\end{definiceBox}
\begin{definiceBox}
\textbf{Vytýkání před závorku} je použití distributivního zákona zprava doleva:
\[AB+AC=A(B+C),\qquad AB-AC=A(B-C).\]
Společný činitel $A$ může být číslo, jednočlen i mnohočlen.
\end{definiceBox}
\poznamka{Rozklad není dělení výrazů. V zápisu $6x^2+9x=3x(2x+3)$ neztrácíme hodnotu $x=0$; rovnici ani zlomek zatím nevytváříme. Vždy ověř zpětným roznásobením.}
\blok{Vytýkání: najdi společný činitel všech členů}
U čísel hledáme společný faktor koeficientů. U proměnných vytýkáme jejich nejnižší společnou mocninu; znaménka zůstávají součástí členů.
\begin{prikladReseniBox}[title={Příklad 1: koeficient a mocnina}]
Rozlož $12x^4-18x^3+6x^2$ na součin.
\tcblower
Společný číselný faktor je $6$, nejnižší mocnina proměnné je $x^2$:
\[12x^4-18x^3+6x^2=6x^2(2x^2-3x+1).\]
Kontrola: $6x^2\cdot 2x^2-6x^2\cdot3x+6x^2=12x^4-18x^3+6x^2$.
\end{prikladReseniBox}
\begin{prikladReseniBox}[title={Příklad 2: více proměnných a mínus}]
Rozlož $-8a^3b^2+12a^2b-4ab$ na součin.
\tcblower
Všemi členy lze dělit $-4ab$, proto
\[-8a^3b^2+12a^2b-4ab=-4ab(2a^2b-3a+1).\]
Zpětným roznásobením ověříme všechny tři členy i znaménka.
\end{prikladReseniBox}
\newpage
\begin{zadaniReseniBox}[title={S1: vytkni největší společný jednočlen}]
\begin{enumerate}[label=\alph*)]
\item $6x^3+9x^2$;
\item $15a^2b-10ab^2+5ab$;
\item $-4x^3+8x^2-12x$;
\item $7y(2y-3)+5(2y-3)$.
\end{enumerate}
\tcblower
\begin{enumerate}[label=\alph*)]
\item $3x^2(2x+3)$;
\item $5ab(3a-2b+1)$;
\item $-4x(x^2-2x+3)$;
\item $(2y-3)(7y+5)$; společnou částí je celá závorka.
\end{enumerate}
\end{zadaniReseniBox}
\blok{Seskupování: společný činitel není na první pohled}
U čtyř členů někdy můžeme vytvořit dvě dvojice tak, aby po vytknutí vznikla \emph{stejná závorka}. Na způsobu seskupení záleží.
\begin{prikladReseniBox}[title={Příklad 3: dva kroky vytýkání}]
Rozlož $ax+ay+bx+by$ na součin.
\tcblower
\begin{align*}
ax+ay+bx+by&=(ax+ay)+(bx+by)\\
&=a(x+y)+b(x+y)\\
&=(a+b)(x+y).
\end{align*}
Zkouška: $(a+b)(x+y)=ax+ay+bx+by$.
\end{prikladReseniBox}
\begin{zadaniReseniBox}[title={S2: seskup a vytkni společnou závorku}]
\begin{enumerate}[label=\alph*)]
\item $x^2+3x+2x+6$;
\item $2ax-6a+bx-3b$;
\item $xy-2x+3y-6$;
\item $x^3+2x^2-x-2$.
\end{enumerate}
\tcblower
\begin{enumerate}[label=\alph*)]
\item $x(x+3)+2(x+3)=(x+2)(x+3)$;
\item $2a(x-3)+b(x-3)=(2a+b)(x-3)$;
\item $x(y-2)+3(y-2)=(x+3)(y-2)$;
\item $x^2(x+2)-(x+2)=(x^2-1)(x+2)$; další rozklad pomocí rozdílu čtverců patří do kombinování metod.
\end{enumerate}
\end{zadaniReseniBox}
\poznamka{U S2 d) zatím postačí $(x^2-1)(x+2)$; rychlí žáci mohou použít vzorec z 46. hodiny a dopsat $(x-1)(x+1)(x+2)$. Není nutné pro všechny postupy vždy dojít ke stejnému počtu činitelů.}
\blok{Volitelná rezerva: poznáš chybu?}
\begin{zadaniReseniBox}[title={R1: zdůvodni správnost}]
Žák napsal $6x^2-9x=3x(2x-9)$. Je rozklad správný? Oprav jej a ověř.
\tcblower
Ne: zpětné roznásobení dává $6x^2-27x$. Správně $6x^2-9x=3x(2x-3)$.
\end{zadaniReseniBox}
\blok{Domácí procvičení}
\begin{zadaniReseniBox}[title={D1: rozlož na součin}]
\begin{enumerate}[label=\alph*)]
\item $8x^3-12x^2$;
\item $9a^3b-6a^2b^2+3ab$;
\item $-10x^2+15x-5$;
\item $3x(x-4)-2(x-4)$;
\item $x^2+5x+2x+10$;
\item $3ay-6a+2by-4b$;
\item $x^3-4x^2+2x-8$;
\item $4x^2y-8xy+6x-12$.
\end{enumerate}
\tcblower
\begin{enumerate}[label=\alph*)]
\item $4x^2(2x-3)$;
\item $3ab(3a^2-2ab+1)$;
\item $-5(2x^2-3x+1)$;
\item $(x-4)(3x-2)$;
\item $(x+2)(x+5)$;
\item $(3a+2b)(y-2)$;
\item $(x-4)(x^2+2)$;
\item $2(x-2)(2xy+3)$.
\end{enumerate}
\end{zadaniReseniBox}
\poznamka{Základ domácí práce a)--f); g)--h) náročnější. Nejprve hledat společný činitel, pak vhodné seskupení. Správnost každý doloží krátkým zpětným roznásobením alespoň u jedné úlohy.}
\end{document}
